\section{Disjoint pair resamplings}
\label{dis:section}

The resamplings indexed by disjoint vertex pairs need not commute:
they both alter the four edges between the two pairs. Carstens, Berger,
and Strona already noted this dependence of undirected
trades~\cite[Section~2.3 of the arXiv version]{CBS18}. We show that
one interaction is bounded below by a cross-edge switch term minus two
errors. Each error is a projection onto an equal-assignment indicator
in one pair fiber. Section~\ref{sl:section} will bound the sum of these
errors by the available switch energy.

For disjoint vertex pairs $a,b$, let $T_{a,b}$ toggle the four edges between
$a$ and $b$ when their $2$-by-$2$ cross matrix has every row and column sum
one, and leave the graph unchanged otherwise. This map is a degree-preserving
involution, and hence acts as a self-adjoint isometry on functions. Define
\begin{equation}
 g_{a,b}=\frac{I-T_{a,b}}2=g_{b,a}.
 \label{sl:def-cross}
\end{equation}
Thus $g_{a,b}$ is an orthogonal projection.

Fix an $a$-fiber $\mathcal F$, and write $b=\{k,l\}$. If both $k$ and $l$
are singleton vertices in this fiber, let $x_k,x_l$ indicate assignment to
the first endpoint of $a$, and set
\begin{equation}
 z_{kl}^{\mathcal F}
 =\mathbf 1_{\{x_k=x_l\}}
    -\mathbb P(x_k=x_l\mid\mathcal F).
 \label{sl:def-equality}
\end{equation}
On $\mathcal F$, define $R_{a,b}$ to be orthogonal projection onto the span
of $z_{kl}^{\mathcal F}$, or zero if this function is zero. If either vertex
is not a singleton, also use the zero operator. Taking the direct sum over
all $a$-fibers defines an orthogonal projection $R_{a,b}$ on the graph
space. Unlike $g_{a,b}$, this projection is directed: $R_{a,b}$ and
$R_{b,a}$ use different fibers. Its definition is independent of the chosen
ordering of the two endpoints of $a$.

\begin{proposition}[Disjoint-pair interaction]\label{dis:interaction}
For disjoint vertex pairs $a,b$ and every real function $f$ on
$\Omega_d$,
\begin{equation}\label{dis:bound}
  2\langle h_a f,h_b f\rangle
  \ge 2\|g_{a,b}f\|^2
       -\|R_{a,b}f\|^2-\|R_{b,a}f\|^2.
\end{equation}
\end{proposition}

\begin{proof}
Write $a=\{i,j\}$ and $b=\{k,l\}$, with $i$ and $k$ the first
endpoints in the orders used to define the slice projections, and put
$O=V\setminus(a\cup b)$. Partition $\Omega_d$ by recording the graph
on $O$, the edges $ij,kl$, and, separately for each $w\in O$, its
numbers of neighbors in $a$ and in $b$. Both pair resamplings preserve
these data. In particular, each full $a$-fiber or $b$-fiber is contained
in a cell of this partition. We prove the inequality on each nonempty
cell, using its conditional uniform measure throughout.

\paragraph{The cell and its two conditional expectations.}
Let $A_0,B_0\subseteq O$ be the outside vertices with exactly one
neighbor in $a,b$, respectively. A graph $G$ in the cell is encoded by
\[
 M=\begin{pmatrix}
       \mathbf 1_{\{ik\in E(G)\}}&\mathbf 1_{\{il\in E(G)\}}\\
       \mathbf 1_{\{jk\in E(G)\}}&\mathbf 1_{\{jl\in E(G)\}}
    \end{pmatrix},
 \qquad U\subseteq A_0,\qquad W\subseteq B_0,
\]
where $U$ contains those outside singletons assigned to $i$, and $W$
those assigned to $k$. The total $s$ of the four entries of $M$ is fixed
by the sum of the degrees in $a$ and the recorded data. Set
\[
 r=M_{ik}+M_{il},\qquad c=M_{ik}+M_{jk}.
\]
For fixed integers $q_0,q_1$, the remaining degree constraints are
exactly
\begin{equation}\label{dis:sizes}
 |U|=q_0-r,\qquad |W|=q_1-c.
\end{equation}
Indeed, these equations enforce the degrees of $i,k$; the degrees of
$j,l$ then follow from the fixed pair totals. Each outside degree is
already enforced by the recorded graph and counts. This also proves
that every binary matrix of total $s$ and every pair of subsets
satisfying \eqref{dis:sizes} give a valid graph in the cell. The sets
$A_0,B_0$ may overlap: the two assignments use different edges and
preserve the two recorded neighbor counts separately. Thus overlap
adds no constraint. All these encodings have the same probability.

Within the cell, the full $a$-fiber data are equivalent to
$X=(c,W)$. Namely, $c$ and $s-c$ specify the counts from $k,l$ to $a$,
and $W$ specifies all remaining edges not incident to $a$. Together
with the cell data this is exactly the defining information of an
$a$-fiber. Similarly, the full $b$-fiber data are $Y=(r,U)$. Therefore
\begin{equation}\label{dis:conditional}
 E_a=\mathbb E[\,\cdot\mid X],\qquad
 E_b=\mathbb E[\,\cdot\mid Y].
\end{equation}
Let $\mathcal A,\mathcal B$ be the subspaces of functions of $X,Y$,
respectively, with constants included, and let
$\mathcal J=(\mathcal A+\mathcal B)^\perp$.
The involution defining $g_{a,b}$ preserves $X$ and $Y$. Hence its
antisymmetric projection kills both subspaces, and, by
self-adjointness,
\begin{equation}\label{dis:g-range}
 \operatorname{ran}g_{a,b}\subseteq\mathcal J.
\end{equation}

\paragraph{Independent cases.}
If $s=0$ or $s=4$, the cross matrix is fixed and $X,Y$ are independent.
If $s=1$ or $s=3$, every feasible combination of its row total $r$ and
column total $c$ determines a unique matrix: it specifies the position
of its unique one or its unique zero. Since the subset restrictions
in \eqref{dis:sizes} are separate, $X,Y$ again have a product law.

For $s=2$, all six cross matrices appear in the following display,
placed according to their first row sum $r$ and first column sum $c$:
\begin{equation}\label{dis:table}
\begin{array}{c|ccc}
 &c=0&c=1&c=2\\ \hline
 r=0&\varnothing&
 \begin{psmallmatrix}0&0\\1&1\end{psmallmatrix}&\varnothing\\[8pt]
 r=1&\begin{psmallmatrix}0&1\\0&1\end{psmallmatrix}&
 \begin{psmallmatrix}1&0\\0&1\end{psmallmatrix},\quad
 \begin{psmallmatrix}0&1\\1&0\end{psmallmatrix}&
 \begin{psmallmatrix}1&0\\1&0\end{psmallmatrix}\\[8pt]
 r=2&\varnothing&
 \begin{psmallmatrix}1&1\\0&0\end{psmallmatrix}&\varnothing
\end{array}.
\end{equation}
Let $m_{rc}$ count the matrices in cell $(r,c)$ of this display:
$m_{11}=2$, the four adjacent cells have count one, and the corners
have count zero. Write $D_r=\{r\ne1\}$ and $D_c=\{c\ne1\}$.
The top and bottom matrices realize $D_r$, and the side matrices
realize $D_c$, so the two events are disjoint. Call $r=c=1$ the center;
its two matrices are exchanged by $T_{a,b}$. If both events have positive
probability, the center is feasible: occurrence of $D_r$ implies
feasibility of a subset $W$ of size $q_1-1$, and occurrence of $D_c$
implies feasibility of a subset $U$ of size $q_0-1$. Together these
subsets realize either central matrix. Thus, if the center is absent,
at most one of the events occurs, and the support lies in a single
row or column of \eqref{dis:table}; the law of $X,Y$ is a product.
If the center is present but one of the events has probability zero,
the same conclusion holds. For example, if $r=1$ always, then $U$ is
uniform of fixed size, independently of $c,W$; the multiplicity in
\eqref{dis:table} depends only on $c$.

Whenever $X,Y$ are independent, both $E_aE_b$ and $E_bE_a$ equal the
projection onto constants in the cell. Consequently $h_a h_b$ is the
orthogonal projection onto $\mathcal J$, and
\[
 2\langle h_a f,h_b f\rangle
   =2\|\operatorname{proj}_{\mathcal J}f\|^2
   \ge 2\|g_{a,b}f\|^2
\]
by \eqref{dis:g-range}. This is stronger than \eqref{dis:bound}.

\paragraph{The remaining correlation has rank one.}
It remains to consider $s=2$ when both events and the center have
positive probability. Put $p_r=\mathbb P(D_r)$ and
$p_c=\mathbb P(D_c)$, and define centered unit vectors
\[
 u=\frac{\mathbf 1_{D_c}-p_c}{\sqrt{p_c(1-p_c)}}\in\mathcal A,
 \qquad
 v=\frac{\mathbf 1_{D_r}-p_r}{\sqrt{p_r(1-p_r)}}\in\mathcal B.
\]
Their correlation is
\begin{equation}\label{dis:rho}
 \langle u,v\rangle=-\rho,\qquad
 \rho=\sqrt{\frac{p_cp_r}{(1-p_c)(1-p_r)}}\in(0,1),
\end{equation}
since $D_r,D_c$ are disjoint and $1-p_c-p_r$ is the positive center
probability.

With the counts $m_{rc}$ associated with \eqref{dis:table}, for every
occurring $X=(c,W)$ and each feasible subset $U$ of size $q_0-r$,
\[
 \mathbb P\bigl(Y=(r,U)\mid X=(c,W)\bigr)
 =\frac{m_{rc}}{\displaystyle\sum_{t=0}^2
       \binom{|A_0|}{q_0-t}m_{tc}},
\]
where a binomial coefficient is zero when its lower index is outside
$\{0,\ldots,|A_0|\}$. Thus the conditional law of $Y$ given $X$
depends only on whether $c=1$. When $c\ne1$, necessarily $r=1$ and
$U$ is uniform of size $q_0-1$; when $c=1$, the displayed formula has
no dependence on $W$.
Symmetrically, the conditional law of $X$ given $Y$ depends only on
whether $r=1$. It follows that
\[
 E_a\mathcal B\subseteq\operatorname{span}(1,u),\qquad
 E_b\mathcal A\subseteq\operatorname{span}(1,v),\qquad
 E_a v=-\rho u,\quad E_b u=-\rho v.
\]
Let $\mathcal A'$ be the centered functions in $\mathcal A$
orthogonal to $u$, and define $\mathcal B'$ symmetrically. The
conditional-expectation inclusions show that
$\mathcal A'\perp\mathcal B$ and $\mathcal B'\perp\mathcal A$.
We therefore have the orthogonal decomposition
\begin{equation}\label{dis:decomposition}
 \operatorname{span}(1)\ \oplus\ \mathcal A'\ \oplus\
 \mathcal B'\ \oplus\ \operatorname{span}(u,v)\ \oplus\ \mathcal J.
\end{equation}
All these summands are invariant under $E_a,E_b$. The first three
contribute zero to $\langle h_a f,h_b f\rangle$, and $\mathcal J$
contributes $\|\operatorname{proj}_{\mathcal J}f\|^2$.

On the plane in \eqref{dis:decomposition}, write the component of
$f$ as $xu+yv$. Its images under $h_a$ and $h_b$, respectively, are
\[
 y(v+\rho u),\qquad x(u+\rho v).
\]
The parenthesized vectors have squared norm $1-\rho^2$ and mutual
inner product $\rho(1-\rho^2)$. Hence the plane contribution obeys
\begin{equation}\label{dis:plane}
 2xy\rho(1-\rho^2)
 \ge -(x^2+y^2)(1-\rho^2)
 =-\|y(v+\rho u)\|^2-\|x(u+\rho v)\|^2.
\end{equation}
Here $2|xy|\le x^2+y^2$ and $0<\rho<1$ justify the inequality
for either sign of $xy$.

\paragraph{The error lines belong to the slice projections.}
The vector $y(v+\rho u)$ is the orthogonal projection of the full
function $f$ onto the line spanned by $v+\rho u$: the spanning vector is
orthogonal to $u$, has inner product $1-\rho^2$ with $v$, and
all other summands of \eqref{dis:decomposition} are orthogonal to
the plane. The symmetric assertion holds for $x(u+\rho v)$.

We claim that the first line belongs to
$\operatorname{ran}R_{a,b}$. Examine
$v+\rho u=v-E_a v$ on each $X$-slice, which, by
\eqref{dis:conditional}, is a full global $a$-fiber. If $c\ne1$,
then $D_r$ is deterministically false and this difference is zero.
If $c=1$, both $k,l$ are singleton vertices for the pair $a$.
Moreover, $D_r$ occurs exactly when their assignments to the two
endpoints of $a$ are equal. Writing these assignments as $x_k,x_l$,
on this fiber we have
\[
 v-E_a v=
 \frac{\mathbf 1_{\{x_k=x_l\}}-
       \mathbb P(x_k=x_l\mid X)}{\sqrt{p_r(1-p_r)}}.
\]
This is a scalar multiple of the centered function defining
$R_{a,b}$, including when that function vanishes. Thus the line
belongs to its fiberwise direct-sum range. Likewise the other line
belongs to $\operatorname{ran}R_{b,a}$. Orthogonal projection onto
a subspace has no larger norm than projection onto a containing
subspace: if $\Pi$ projects onto a line contained in
$\operatorname{ran}R$, then $\Pi R=\Pi$ and
$\|Rf\|^2\ge\|\Pi Rf\|^2=\|\Pi f\|^2$. Consequently
\[
 \|y(v+\rho u)\|^2\le\|R_{a,b}f\|^2,\qquad
 \|x(u+\rho v)\|^2\le\|R_{b,a}f\|^2.
\]
Combining these inequalities with \eqref{dis:plane} and
\eqref{dis:g-range} proves \eqref{dis:bound} on the cell.

Finally, $E_a,E_b,g_{a,b}$ preserve the cells. So do $R_{a,b}$ and
$R_{b,a}$, since their defining full pair fibers are contained in
the cells. Averaging the conditional inequalities proves the
asserted global inequality.
\end{proof}

We retain the positive term in Proposition~\ref{dis:interaction} to
control the two directed error terms after summation. The estimate needed for that step is, for each fixed $a$,
\[
 \sum_{b:\,b\cap a=\varnothing}\|R_{a,b}f\|^2
 \le \sum_{b:\,b\cap a=\varnothing}\|g_{a,b}f\|^2.
\]
We prove this estimate next, using the uniform subset structure of
each $a$-fiber.
